\section{Exporting}
\label{sec:export}

Both the main window and the Strain window have an \textbf{Export
Results} button. Both open the \emph{same} dialog. The difference is
only which visualization preset (colormap, range, deformed toggle)
seeds the dialog --- whichever window you opened it from.

\subsection{The Export dialog}

Four tabs at the top of the dialog:

\subsubsection*{Data tab}

Check boxes for numeric-data formats:
\begin{itemize}
  \item \textbf{NPZ} --- NumPy compressed archive; preferred for
        Python downstream.
  \item \textbf{MAT} --- MATLAB v7.3 HDF5 file.
  \item \textbf{CSV} --- one row per node per frame; plain text,
        large.
\end{itemize}

Plus a multi-select list of which fields to include (displacement,
strain components, derived strains, rotation).

\subsubsection*{Images tab}

Per-field list with Enable / Colormap / Range / Background options
for rendering each field as PNG / JPG / TIFF / PDF. One file per
frame per enabled field.

\begin{tipbox}
\textbf{Include colorbar} is ON by default. Images without a
colorbar show coloured patches but no scale, which is almost never
useful for quantitative reports. Uncheck it only if you plan to
composite a manually-placed colorbar in Illustrator / Inkscape.
\end{tipbox}

Adjustable settings:
\begin{description}[leftmargin=1.2cm,style=nextline]
  \item[Image format]
    png / jpg / tif / pdf. PNG is lossless and near-universal; use
    PDF for figures you will embed in LaTeX.
  \item[DPI] 72 / 150 / 300 / 600. 300 is print-quality.
  \item[Background]
    \emph{Reference frame}: all frames rendered on the reference.
    \emph{Current frame}: each frame rendered on its deformed image.
\end{description}

\subsubsection*{Animations tab}

Render an MP4 / GIF for each selected field, sweeping through
frames. Configurable FPS (usually match the frame rate of the
experiment). Colorbar on by default here too.

\begin{warnbox}
MP4 export requires \texttt{ffmpeg} on your system PATH. If ffmpeg is
missing, the dialog will switch silently to GIF and log a warning.
\end{warnbox}

\subsubsection*{Report tab}

Generates a multi-page PDF summarizing the run: parameters, key
fields at mid and last frame, RMSE (if ground truth is available),
and optional thumbnail animation.

\subsection{Destination folder}

At the top of the dialog. Defaults to
\texttt{<image\_folder>/al-dic-export/<timestamp>/}. Click
\emph{Browse} to pick a different location. Exports are written
inside that folder with a user-supplied prefix (default
\texttt{result}).

\subsection{Running an export}

Click the primary \textbf{Export} button at the bottom. Progress
shows per-file (data) and per-frame (images / animations). The
Cancel button stops cleanly; already-written files are kept.

When the export completes, a green success message tells you how
many files were written and their location. The dialog stays open
so you can export more formats; click Close when done.
